%==============================================================================
% FICHAMENTO CRÍTICO — Deep Learning for Automated Detection of Maxillary
% Sinus Lesions on Panoramic Radiographs
% Conforme NBR 6023 (referências) e NBR 10520 (citações)
%==============================================================================
\section{Fichamento: \citeonline{kuwana2021maxillary}}

% --- 1. IDENTIFICAÇÃO ---
\subsection{Identificação}
\begin{tabular}{ll}
\textbf{Título:} & Performance of deep learning object detection technology in the
detection and diagnosis of maxillary sinus lesions on panoramic radiographs \\
\textbf{Autor(es):} & Kuwana, Ryosuke; Ariji, Yoshiko; Fukuda, Motoki; Kise,
Yoshitaka; Nozawa, Michihito; Kuwada, Chiaki; Muramatsu, Chisako; Katsumata,
Akitoshi; Fujita, Hiroshi; Ariji, Eiichiro \\
\textbf{Periódico:} & Dentomaxillofacial Radiology \\
\textbf{Ano:} & 2021 \\
\textbf{Volume:} & 50 \\
\textbf{Número:} & 1 \\
\textbf{Páginas:} & 20200171 \\
\textbf{DOI:} & \url{https://doi.org/10.1259/dmfr.20200171} \\
\textbf{Qualis:} & A2 (Odontologia) \\
\textbf{Tipo:} & Artigo original \\
\end{tabular}

% --- 2. OBJETIVO ---
\subsection{Objetivo}
Determinar o desempenho de uma técnica de deep learning baseada em detecção de
objetos (DetectNet) para detectar seios maxilares em radiografias panorâmicas e
classificar lesões do seio maxilar (sinusite, cistos) em comparação com seios
saudáveis. O estudo visa fornecer uma ferramenta automatizada de auxílio
diagnóstico para cirurgiões-dentistas na interpretação de alterações
sinusais em exames panorâmicos de rotina.

% --- 3. METODOLOGIA ---
\subsection{Metodologia}
\textbf{Desenho do estudo:} Estudo experimental retrospectivo de desenvolvimento
e validação de modelo de deep learning para detecção e classificação de objetos. \\
\textbf{Amostra:} 1.174 imagens de seios maxilares (587 saudáveis, 416 com
inflamação/sinusite, 171 com cistos) extraídas de radiografias panorâmicas,
divididas em conjuntos de treino, teste 1 e teste 2. \\
\textbf{Métricas principais:} Sensibilidade de detecção, taxa de falsos-positivos
por imagem, acurácia, sensibilidade e especificidade para classificação
diagnóstica. \\
\textbf{Validação:} Treinamento com 1.000 épocas usando DetectNet (NVIDIA DIGITS)
com validação em dois conjuntos de teste independentes (teste 1 e teste 2).

% --- 4. RESULTADOS PRINCIPAIS ---
\subsection{Resultados Principais}
\begin{itemize}
  \item Sensibilidade de detecção de 100\% para seios saudáveis e inflamados
        em ambos os conjuntos de teste; para cistos, 98\% (teste 1) e 89\%
        (teste 2).
  \item Taxa de falsos-positivos por imagem de 0,00 para seios saudáveis e
        inflamados; para cistos, 0,00 (teste 1) e 0,18 (teste 2).
  \item Acurácia, sensibilidade e especificidade para sinusite de 90--91\%,
        88--85\% e 91--96\%; para cistos, 97--100\%, 80--100\% e 100--100\%,
        respectivamente.
\end{itemize}

% --- 5. LIMITAÇÕES ---
\subsection{Limitações}
\begin{itemize}
  \item Amostra proveniente de um único centro odontológico japonês, o que
        limita a validade externa e a generalização para outras populações e
        equipamentos de imagem.
  \\item A classificação dos cistos apresentou sensibilidade significativamente
        menor (80--100\%) em comparação com a sinusite, possivelmente devido
        à heterogeneidade morfológica dos cistos no seio maxilar.
  \item Foram excluídas lesões malignas do seio maxilar devido ao pequeno
        número de casos, deixando uma lacuna importante para a aplicação
        clínica completa do modelo.
  \item O estudo não realizou validação externa com equipamentos de diferentes
        fabricantes ou com diferentes protocolos de aquisição de imagem.
\end{itemize}

% --- 6. CONTRIBUIÇÃO PARA O LIVRO ---
\subsection{Contribuição para o Livro}
Este artigo fundamenta diretamente o Capítulo 8, que trata do uso de deep
learning para diagnóstico por imagem em radiografias panorâmicas, com foco
na detecção automatizada de alterações nos seios maxilares. Os resultados
demonstram que modelos de detecção de objetos (DetectNet) podem atingir
sensibilidade de 100\% na identificação de seios maxilares saudáveis e
inflamados, com taxas de falso-positivo próximas de zero. O estudo é
particularmente relevante para a subseção que discute a aplicação de DL
em contexto de anatomia sinusal para planejamento de implantes e cirurgias
ortognáticas. O mapeamento SDD do livro incorpora as métricas de acurácia
($\geq$ 90\%) como critérios de aceitação para modelos de classificação
de lesões sinusais em panorâmicas.

% --- 7. ANÁLISE CRÍTICA ---
\subsection{Análise Crítica}
\begin{itemize}
  \item \textbf{Validade interna:} Alta -- O estudo utiliza dois conjuntos de
        teste independentes (teste 1 e teste 2) para validação, com métricas
        consistentes entre ambos.
  \item \textbf{Validade externa:} Média -- Dados de um único centro e de uma
        única população asiática, sem validação multicêntrica ou em populações
        ocidentais.
  \item \textbf{Reprodutibilidade:} Alta -- O pipeline utilizando DetectNet
        com 1.000 épocas e o ambiente DIGITS é reproduzível; a descrição do
        dataset é detalhada.
  \item \textbf{Relevância clínica:} Alta -- A detecção automatizada de
        sinusite (AUC $\geq$ 0,90) pode auxiliar cirurgiões-dentistas gerais
        na triagem de pacientes que necessitam de encaminhamento
        otorrinolaringológico.
  \item \textbf{Conflito de interesses:} Não -- Os autores declararam ausência
        de conflitos de interesse.
  \item \textbf{Viés identificado:} Viés de seleção -- apenas lesões
        benignas (sinusite e cistos) foram incluídas, excluindo neoplasias
        malignas e outras patologias sinusais relevantes na prática clínica.
\end{itemize}

% --- 8. CITAÇÕES RELEVANTES ---
\subsection{Citações Relevantes}
\begin{citacao}
``Detection sensitivities of healthy (Class 0) and inflamed (Class 1) maxillary
sinuses were 100\% for both testing 1 and testing 2 data sets, whereas they
were 98 and 89\% for cysts of the maxillary sinus regions (Class 2).'' (p. 5).
\end{citacao}

\begin{citacao}
``Accuracies, sensitivities and specificities for diagnosis maxillary sinusitis
were 90--91\%, 88--85\%, and 91--96\%, respectively; for cysts of the maxillary
sinus regions, these values were 97--100\%, 80--100\%, and 100--100\%,
respectively.'' (p. 6).
\end{citacao}

% --- 9. MAPEAMENTO SDD ---
\subsection{Mapeamento SDD}
\textbf{Problema:} Detecção e classificação automatizada de lesões do seio
maxilar (sinusite e cistos) em radiografias panorâmicas, reduzindo a
necessidade de exames complementares (CBCT) e falso-diagnósticos. \\
\textbf{Entrada:} Imagens de radiografias panorâmicas (anonimizadas) com
anotação manual das regiões de interesse (seios maxilares). \\
\textbf{Saída:} Caixas delimitadoras com classificação em três categorias:
saudável, inflamado (sinusite) e cisto. \\
\textbf{Critérios:} Sensibilidade de detecção $\geq$ 98\%; acurácia de
classificação $\geq$ 90\% para sinusite e $\geq$ 95\% para cistos;
taxa de falsos-positivos $<$ 0,05 por imagem. \\
\textbf{Confiança:} 0,78 -- Estudo bem executado, com validação em dois
testes independentes, mas limitado por amostra unicêntrica e exclusão de
lesões malignas.
\end{tabular}

% --- 10. REFERÊNCIA ABNT ---
\subsection{Referência ABNT}
\begin{verbatim}
KUWANA, Ryosuke et al. Performance of deep learning object detection technology
in the detection and diagnosis of maxillary sinus lesions on panoramic
radiographs. Dentomaxillofacial Radiology, London, v. 50, n. 1, p. 20200171,
2021. DOI: 10.1259/dmfr.20200171.
\end{verbatim}
